\chapter{Hỏi Để Làm Việc Nhóm}
\chapterintro{Câu hỏi đúng giúp nhóm bớt ngồi chung mà mạnh ai nấy làm}{Muốn nhóm có thật, phải có câu hỏi buộc các em cần nhau.}

\section{Nếu nhóm chỉ được giữ lại một ý, ý nào sống sót?}
\begin{cauhoibox}{Câu hỏi}
Nếu nhóm chỉ được giữ lại một ý, ý nào sống sót?
\end{cauhoibox}
\begin{taisaocoich}
Câu hỏi này buộc nhóm phải lọc và thống nhất, thay vì ai cũng nói phần mình rồi thôi.
\end{taisaocoich}
\begin{goiybien}
Rất hợp để chốt một vòng thảo luận nhóm ngắn.
\end{goiybien}
\ngancach

\section{Ai trong nhóm nên nói trước và vì sao?}
\begin{cauhoibox}{Câu hỏi}
Ai trong nhóm nên nói trước và vì sao?
\end{cauhoibox}
\begin{taisaocoich}
Câu này làm nhóm phải nhìn vào vai trò chứ không chỉ nhìn vào nội dung. Nó giúp tránh cảnh một bạn ôm hết phần nói.
\end{taisaocoich}
\begin{goiybien}
Dùng rất tốt với lớp còn yếu kỹ năng nhóm hoặc đang có nhóm trưởng áp đảo.
\end{goiybien}
\ngancach

\section{Nhóm mình đang thiếu ý hay thiếu nghe nhau?}
\begin{cauhoibox}{Câu hỏi}
Nhóm mình đang thiếu ý hay thiếu nghe nhau?
\end{cauhoibox}
\begin{taisaocoich}
Đây là câu hỏi sửa nhịp nhóm rất nhanh. Nhiều khi vấn đề không phải nghèo ý tưởng, mà là không lắng nghe nhau đủ.
\end{taisaocoich}
\begin{goiybien}
Cho nhóm trả lời bằng giơ tay hoặc chọn một từ, không cần nói dài.
\end{goiybien}
\ngancach

\section{Nhóm mình đang trùng ý hay trùng tiếng?}
\begin{cauhoibox}{Câu hỏi}
Nhóm mình đang trùng ý hay trùng tiếng?
\end{cauhoibox}
\begin{taisaocoich}
Câu hỏi này giúp nhóm tự nhận ra vấn đề rất nhanh. Nhiều khi nhóm tưởng đang làm việc sôi nổi, nhưng thật ra chỉ đang ồn lên.
\end{taisaocoich}
\begin{goiybien}
Hợp với những lúc thảo luận nhóm bắt đầu loãng hoặc ai cũng muốn chen phần mình.
\end{goiybien}
\ngancach

\section{Nếu đổi vai trong nhóm, ai nên thử vai nào?}
\begin{cauhoibox}{Câu hỏi}
Nếu đổi vai trong nhóm một lượt, ai nên thử vai nào?
\end{cauhoibox}
\begin{taisaocoich}
Đổi vai làm học sinh nhìn lại cách nhóm đang vận hành. Nó đặc biệt tốt khi vài bạn nói quá nhiều còn vài bạn gần như mất vai trò.
\end{taisaocoich}
\begin{goiybien}
Có thể chỉ đổi thử trong 3 phút: người trình bày thành người ghi ý, người im hơn thành người chốt câu cuối.
\end{goiybien}
